\section{Coefficient spaces and clamped Poisson estimates}\label{inf:sec-spaces}

From now on $p\in\tfrac12\Z$ and the integer split are fixed. Constants in operator
estimates are uniform over $p\ge64$ and admissible splits unless a
subscript explicitly permits dependence on the centre order.
The angular coefficient radius is fixed at $\rho_w=2$ throughout
the Newton argument. Larger radii will be used only to estimate the
known centre.

\subsection{The source, shape, and clamped subspaces}

Define the degree weight
\begin{equation}\label{inf:degree-weight}
 \wt(\ell)=(1+\ell/p^2)^p,\qquad \ell\ge0.
\end{equation}
It is increasing and obeys
\begin{equation}\label{inf:weight-properties}
 \wt(\ell_1+\ell_2)\le\wt(\ell_1)\wt(\ell_2),\qquad
 \frac{\wt(\ell+2)}{\wt(\ell)}\le e^{2/p},\qquad
 \wt(\ell)\le e^{\ell/p}.
\end{equation}
These follow from $1+u+v\le(1+u)(1+v)$ and $1+u\le e^u$.

\begin{definition}[The fixed Banach spaces]\label{inf:spaces}
For the polynomial basis \eqref{inf:radial-basis}, let $\Cc=\Cc_{p,2}$
be the completion with norm
\begin{equation}\label{inf:C-norm}
 \nrm f_\Cc=\sum_{j,s\ge0}2^j\wt(2j+2s)|f_{j,s}|.
\end{equation}
Its moment spaces have norms
\begin{equation}\label{inf:C-moments}
 \nrm f_{\Cc^{[k]}}=
 \sum_{j,s\ge0}2^j\wt(2j+2s)(1+2j+2s)^k|f_{j,s}|.
\end{equation}
Let $\Gc_p=\{g\in\Cc:g_{j,0}=0\text{ for every }j\}$ and let
$\Ec_1$ be the space of all boundary coefficient sequences with
finite norm
\begin{equation}\label{inf:E-norm}
 h=\sum_{j\ge0}h_jG_j,\qquad
 \nrm h_{\Ec_1}=\sum_{j\ge0}(1+j)2^j\wt(2j)|h_j|.
\end{equation}
Write $\HP  h=\sum_jh_jS_j$ for its harmonic extension, and
$\Xc_p=\Gc_p\times\Ec_1$ with the sum norm.
The Newton spaces are real; their complexifications are used for
analytic maps and estimates. Dot products in algebraic differential
identities are extended complex-bilinearly; Hilbert inner products
and norms retain their Hermitian meaning.
\end{definition}

\begin{lemma}[Completions and shape regularity]\label{inf:space-embedding}
The space $\Cc$ embeds injectively into invariant normalized
$L^2(B)$ with norm at most one. Polynomial partial sums are dense.
For each fixed $k$, harmonic extension sends $\Ec_1$ continuously
into $C^k(\overline B)$, with constants depending on $k$ and the
dimension. The first two boundary derivative bounds are independent
of $p$.
In particular
\begin{equation}\label{inf:shape-C2}
 \nrm h_{C^2(S^{2p-1})}\le C\nrm h_{\Ec_1},\qquad
 \nrm{D((\HP h)y)}_{\mathrm{op},\infty}\le C\nrm h_{\Ec_1}.
\end{equation}
The boundary function $h$ is real analytic.
\end{lemma}
\begin{proof}
By \eqref{inf:radial-norm}, every basis function has $L^2$ norm at
most one. The triangle inequality gives the continuous embedding
on finite sums and hence on their completion. Its injectivity
follows by taking the inner product with each orthogonal basis
function: all coefficients of a zero $L^2$ limit vanish.
Weighted $\ell^1$ completeness and truncation give the remaining
Banach-space assertions.

Lemma~\ref{inf:angular-sup} bounds the first two boundary
derivatives by $2j$ and $4j^2$. The ratios of these quantities to
$(1+j)2^j\wt(2j)$ are uniformly bounded. Summing proves the first
estimate in \eqref{inf:shape-C2}. Radial homogeneity gives
$\partial_\varrho S_j=2j\varrho^{2j-1}G_j$;
the tangential bound gives $|\nabla S_j|\le2\sqrt2j$ in the ball.
The product rule proves the second estimate.
For the full extension, choose any fixed $1<R_*<\sqrt2$.
Homogeneity and the endpoint bound give
$|S_j(y)|\le R_*^{2j}$ for real $|y|\le R_*$.
Thus the series for $\HP h$ converges uniformly on each compact
subset of $B_{\sqrt2}$, with supremum there controlled by the
shape norm. Its finite sums are harmonic polynomials. Passing
to the limit in the mean-value identity shows that the limit is
harmonic in that larger ball. The Poisson formula on a ball of
fixed radius about each point of $\overline B$ can therefore be
differentiated any fixed number of times. It gives uniform
convergence of all those derivatives on $\overline B$, with
constants depending on the derivative order and fixed dimension.
It also proves real analyticity of the extension on a
neighbourhood of $\overline B$. The dimension-independent first
two boundary bounds are the separate estimates above.

To see analyticity of the boundary directly, expand $G_j$ in
Chebyshev polynomials on $[0,1]$. Its real supremum is at most one,
so each coefficient has absolute value at most two. On any
Bernstein ellipse of parameter $1<\rho_{\mathrm{ell}}<2$ this gives
$|G_j(z)|\le2(j+1)\rho_{\mathrm{ell}}^j$.
The series for $h$ converges uniformly on that ellipse, since
$\rho_{\mathrm{ell}}<2$, and defines a holomorphic function there. Composing
with the polynomial $x=|X|^2$ on the real sphere proves real
analyticity of the boundary function.
\end{proof}

\subsection{The positive multiplier module}

For endpoint-normalized Jacobi polynomials with parameters
$\alpha,\beta>-1$, put $u=\alpha+\beta+1$ and $v=\alpha-\beta$.
Gasper's theorem \cite[Theorem~1]{inf:gasper} gives nonnegative
linearization coefficients provided
\begin{equation}\label{inf:gasper-hypothesis}
 \alpha\ge\beta,\qquad
 u(u+5)(u+3)^2\ge(u^2-7u-24)v^2.
\end{equation}
The normalization is at the endpoint corresponding to the larger
Jacobi exponent $\alpha$. In the two-block parameters
$\ab=a/2\le\bb=p-a/2$, this is $x=1$, the endpoint of the
factor $(1-x)^{\bb-1}$.
Here $\alpha=\bb-1\ge\beta=\ab-1$, $u=p-1$, and
$0\le v=\bb-\ab\le p/2$. For $p\ge64$ the margin in
\eqref{inf:gasper-hypothesis} is at least
\[
 \frac{3p^4+37p^3+64p^2-16p-64}{4}>0.
\]
Thus in our normalization
\begin{equation}\label{inf:gasper-product}
 G_iG_j=\sum_{|i-j|\le k\le i+j}c_{ij}^kG_k,
 \qquad c_{ij}^k\ge0,\qquad\sum_kc_{ij}^k=1.
\end{equation}
The upper support follows from polynomial degree; the lower support
follows from orthogonality, since a polynomial of degree below $i$
has zero inner product with $G_i$. The sum equals one by evaluating
at $x=1$.

\begin{lemma}[Positive radial stencils]\label{inf:radial-stencils}
For $\beta\ge0$ and $s\ge0$,
\begin{align}
 J_s^\beta&=\frac{s+\beta+1}{2s+\beta+1}J_s^{\beta+1}
              +\frac{s}{2s+\beta+1}J_{s-1}^{\beta+1},
 \label{inf:raising-stencil}\\
 tJ_s^{\beta+1}&=\frac{s+\beta+1}{2s+\beta+2}J_s^\beta
              +\frac{s+1}{2s+\beta+2}J_{s+1}^\beta.
 \label{inf:lowering-stencil}
\end{align}
The term with index $-1$ in the first identity is zero.
All coefficients are nonnegative and sum to one.
Multiplication by $t$ in a fixed radial basis also has three
nonnegative coefficients of sum one.
\end{lemma}
\begin{proof}
Insert the finite binomial formula
\eqref{inf:radial-expansion}. Comparing each power of $t$ and using
$\binom{z+1}{s}=\binom z s+\binom z{s-1}$ gives
\eqref{inf:raising-stencil} and \eqref{inf:lowering-stencil}.
Their coefficients have the displayed signs and sums.
For multiplication by $t$ in a fixed basis, orthogonality makes
its matrix tridiagonal. The upper coefficient is the positive
ratio of leading coefficients; symmetry and the positive squared
norms give a positive lower coefficient. The diagonal coefficient
is $\int t(J_s^\beta)^2t^\beta dt/\int(J_s^\beta)^2t^\beta dt$,
which is nonnegative. Evaluation at $t=1$ gives their sum one.
\end{proof}

\begin{proposition}[Solid multipliers and the shape algebra]
\label{inf:module}
For integers $i,h\ge0$, multiplication by $S_it^h$ is bounded on $\Cc$ and
\begin{equation}\label{inf:module-bound}
 \nrm{M_{S_it^h}}_{\Cc\to\Cc}\le2^i\wt(2i+2h).
\end{equation}
The boundary space $\Ec_1$ is a Banach algebra, with
$\nrm{h_1h_2}_{\Ec_1}\le\nrm{h_1}_{\Ec_1}\nrm{h_2}_{\Ec_1}$.
In particular
\begin{equation}\label{inf:harmonic-module}
 \nrm{M_{\HP h}}\le\nrm h_{\Ec_1},\qquad
 \nrm{M_{\Euler\HP h}}\le2\nrm h_{\Ec_1}.
\end{equation}
Here $\Euler=y\cdot\nabla$.
\end{proposition}
\begin{proof}
A product $S_it^h\Psi_{j,s}$ has, by
\eqref{inf:gasper-product}, output profiles
$t^{i+j-k+h}J_s^{\beta_j}$ with positive angular coefficients of
sum one. If $k\ge j$, raise the radial parameter by $2(k-j)$
using \eqref{inf:raising-stencil}, then multiply by the remaining
powers of $t$. If $k<j$, use $2(j-k)$ powers for lowering by
\eqref{inf:lowering-stencil}. There are enough powers because
$i+j-k\ge2(j-k)$ is equivalent to $k\ge j-i$, part of the
triangle support. The remaining powers use the fixed-parameter
positive multiplication stencil.
Every conversion has coefficient sum one. The output total degree
is at most the input degree plus $2i+2h$, and $2^k\le2^{i+j}$.
Equation~\eqref{inf:weight-properties} proves
\eqref{inf:module-bound}, first on basis columns and then by
absolute summation and completion.

For boundary products use the same angular support,
$1+k\le(1+i)(1+j)$, and
$\wt(2k)\le\wt(2i)\wt(2j)$.
This proves the algebra inequality. The harmonic multiplier
bounds follow by summing \eqref{inf:module-bound} and using
$\Euler S_i=2iS_i$.
These coefficient multipliers equal pointwise multiplication
on $\Cc$: equality holds on finite source sums, which converge
in $\Cc$ and hence in $L^2$, while multiplication by a fixed
solid polynomial is continuous in $L^2$.
For a shape extension its finite sums converge uniformly on
the ball and in multiplier norm, so both limits give the same
pointwise multiplier.
\end{proof}

\subsection{Exact Poisson columns and traces}

Let $\KD=\Delta_D^{-1}$ on the unit ball, so $\KD=-H_0^{-1}$.
Its sign is part of this convention. The same operator restricted
to $\Gc_p$ will occasionally be denoted by $K$.

\begin{lemma}[Poisson columns]\label{inf:poisson-columns}
Fix an angular input $j$, and put $B=2j$, $\beta=B+p-1$,
$b_0=\beta+1$, and $d_s=\beta+2s$.
For $s\ge1$,
\begin{align}
 \KD\Psi_{j,s}&=S_jk_s(t),\nonumber\\
 k_s&=\frac{J_{s-1}^\beta}{4d_s(d_s+1)}
 -\frac{J_s^\beta}{2d_s(d_s+2)}
 +\frac{J_{s+1}^\beta}{4(d_s+1)(d_s+2)}.\label{inf:poisson-profile}
\end{align}
For $s=0$,
\begin{equation}\label{inf:poisson-harmonic}
 k_0=\frac{t-1}{4b_0}
     =\frac{J_1^\beta-J_0^\beta}{4b_0(b_0+1)}.
\end{equation}
For $s\ge1$ the derivative identities are
\begin{align}
 k_s'&=\frac{J_s^{\beta+1}-J_{s-1}^{\beta+1}}{4(d_s+1)},
 \label{inf:poisson-derivative}\\
 \beta k_s+tk_s'&=
 \frac{J_{s+1}^{\beta-1}-J_s^{\beta-1}}{4(d_s+1)}.
 \label{inf:poisson-lowered}
\end{align}
For $s=0$ their replacements are
$k_0'=1/(4b_0)$ and
$\beta k_0+tk_0'=J_1^{\beta-1}/(4b_0)$.
\end{lemma}
\begin{proof}
Write $J_s^\beta(t)=\sum_{m=0}^s a_{s,m}t^m$, where
\eqref{inf:radial-expansion} gives
\[
 a_{s,m}=(-1)^{s-m}\binom sm\binom{s+\beta+m}s.
\]
The coefficient equation for
$4\{tk_s''+(\beta+1)k_s'\}=J_s^\beta$ is
\[
 [t^{m+1}]k_s=\frac{a_{s,m}}{4(m+1)(\beta+m+1)}.
\]
For $1\le m\le s$, the adjacent-layer ratios are
\[
 \frac{a_{s-1,m}}{a_{s,m}}=-\frac{s-m}{s+\beta+m},\qquad
 \frac{a_{s+1,m}}{a_{s,m}}=-\frac{s+\beta+m+1}{s+1-m},
\]
whereas
\[
 \frac{a_{s,m-1}}{a_{s,m}}
 =-\frac{m(\beta+m)}{(s+1-m)(s+\beta+m)}.
\]
With $d=d_s=\beta+2s$, the coefficient in the proposed formula,
divided by $a_{s,m}$, is therefore
\begin{align*}
 &-\frac{s-m}{4d(d+1)(s+\beta+m)}-\frac1{2d(d+2)}
 -\frac{s+\beta+m+1}{4(d+1)(d+2)(s+1-m)}\\
 &\hspace{35mm}=-\frac1{4(s+1-m)(s+\beta+m)}.
\end{align*}
The equality follows by multiplication by the common denominator
and substitution of $d=\beta+2s$; its right side is exactly
$a_{s,m-1}/\{4m(\beta+m)a_{s,m}\}$.
For the top coefficient use
\[
 \frac{a_{s+1,s+1}}{a_{s,s}}
 =\frac{(d+2)(d+1)}{(s+1)(s+\beta+1)}.
\]
Finally the three column coefficients sum to zero, fixing the
constant coefficient by $k_s(1)=0$. This proves
\eqref{inf:poisson-profile} for every $s\ge1$.

For \eqref{inf:poisson-derivative}, the coefficient of $t^m$ in
$J_s^{\beta+1}-J_{s-1}^{\beta+1}$, divided by $a_{s,m}$, equals
\[
 \frac{s+\beta+m+1}{\beta+m+1}
 +\frac{s-m}{\beta+m+1}
 =\frac{d+1}{\beta+m+1}.
\]
This is exactly $4(d+1)[t^m]k_s'/a_{s,m}$.
For \eqref{inf:poisson-lowered}, at $m\ge1$ the left coefficient is
$(\beta+m)[t^m]k_s=a_{s,m-1}/(4m)$.
The two binomial ratios in the right numerator give
\[
 \frac{[t^m](J_{s+1}^{\beta-1}-J_s^{\beta-1})}{a_{s,m-1}}
 =\frac{s+\beta+m}{m}+\frac{s+1-m}{m}
 =\frac{d+1}{m}.
\]
Their nonconstant coefficients agree. The previous derivative
identity gives $k_s'(1)=0$; both sides of
\eqref{inf:poisson-lowered} vanish at one, so their constants
agree as well.
For $s=0$, the elementary expression
\eqref{inf:poisson-harmonic} solves the equation and gives both
replacement derivative identities directly. In particular this
layer has the denominator $4b_0$ in \eqref{inf:poisson-harmonic}.
By \eqref{inf:profile-laplacian}, these polynomial functions solve
the Dirichlet Poisson problem. Uniqueness identifies them with
$\KD\Psi_{j,s}$.
\end{proof}

\begin{proposition}[Poisson bounds and exact clamping]
\label{inf:poisson-bounds}
On all of $\Cc$, including the layer $s=0$,
\begin{equation}\label{inf:poisson-operator-bounds}
 \nrm{\KD}\le\frac3{p(p+1)},\qquad
 \nrm{\Euler\KD}\le C/p,\qquad
 \nrm{\Euler^2\KD}\le C.
\end{equation}
More precisely, its weighted basis columns satisfy, for all $j,s\ge0$,
\begin{align}
 \nrm{\KD\Psi_{j,s}}_\Cc
 &\le\frac{C}{(p+2j+2s)^2}\nrm{\Psi_{j,s}}_\Cc,\nonumber\\
 \nrm{\Euler\KD\Psi_{j,s}}_\Cc
 &\le\frac{C}{p-1+2j+2s}\nrm{\Psi_{j,s}}_\Cc.
 \label{inf:poisson-column-bounds}
\end{align}
The normal trace is the bounded map
\begin{equation}\label{inf:poisson-trace}
 \partial_\nu\KD f=
 \sum_{j\ge0}\frac{f_{j,0}}{2(2j+p)}G_j,
 \qquad
 \nrm{\partial_\nu\KD}_{\Cc\to\Ec_1}\le\tfrac12.
\end{equation}
Consequently $U=1+\KD g$ has Dirichlet trace one and normal trace
zero exactly when $g\in\Gc_p$.
\end{proposition}
\begin{proof}
The absolute sums of the columns in
\eqref{inf:poisson-profile} and \eqref{inf:poisson-harmonic} are
$1/(d_s(d_s+2))$ and $1/(2b_0(b_0+1))$.
The only upward radial shift increases degree by two, whose
weight ratio is at most $e^{2/p}<2$. This proves the first bound.
On one profile,
\[
 \Euler\KD=Bk_s+2tk'_s,\qquad
 \Euler^2\KD=B^2k_s+4(1-p)tk'_s+tJ_s^\beta.
\]
The second equality uses the Poisson equation to eliminate $k_s''$.
The derivative column in parameter $\beta+1$ has sum at most
$1/(2(d_s+1))$; multiplication by $t$ lowers it positively to
parameter $\beta$. The separate $s=0$ column gives the same bound
up to an absolute constant. Since $B\le d_s$ and $d_s\ge p-1$,
these formulas prove the other two operator bounds and
\eqref{inf:poisson-column-bounds}.

At the boundary $S_j=G_j$ and
$\partial_\nu(S_jk_s)=2jk_s(1)G_j+2k'_s(1)G_j$.
Only $s=0$ contributes, giving \eqref{inf:poisson-trace}.
Its weighted trace bound follows from
$(1+j)/(2(2j+p))\le1/2$ and
$\wt(2j)\le\wt(2j+2s)$.
For a general coefficient source, approximate by finite sums.
The embedding in $L^2$ and Dirichlet $H^2$ regularity give
$\KD f_n\to\KD f$ in $H^2(B)$.
The displayed trace maps converge in $\Ec_1$ and thus in every
fixed boundary Sobolev norm, by exponential angular decay.
The Sobolev trace theorem identifies this limit with the actual
normal trace of the $H^2$ solution. Dirichlet traces are zero by
Poisson construction. Finally orthogonality of the $G_j$ shows
that the trace vanishes if and only if each $f_{j,0}$ vanishes.
This proves the asserted characterization on the completed spaces.
\end{proof}

\subsection{Raising coupling and the clamped gain}

\begin{lemma}[Contraction under parameter raising]\label{inf:raising-coupling}
For $\beta\ge0$ and integers $s\ge1$, $m\ge0$, the coefficient
vectors of $J_s^\beta$ and $J_{s-1}^\beta$ in parameter
$\beta+m$ are probability vectors: their entries are nonnegative
and sum to one by \eqref{inf:raising-stencil}. Their $\ell^1$
distance is at most
\begin{equation}\label{inf:coupling-bound}
 2e^2\frac z{z+m},\qquad z=2s+\beta+1.
\end{equation}
The bound is uniform in all its parameters.
\end{lemma}
\begin{proof}
One raising step in \eqref{inf:raising-stencil} keeps index
$\sigma$ or lowers it by one, with down probability
$q_\sigma(b)=\sigma/(2\sigma+b+1)$.
Couple the chains started at $s,s-1$ with one uniform random
number at each step. Since $q_\sigma$ increases in $\sigma$,
the indices remain adjacent or merge, and after merging remain
equal. Before merging, at step $h$ their upper index is at most
$s$ and their parameter is $b=\beta+h$. Their merging probability
is at least
\[
 q_\sigma(b)-q_{\sigma-1}(b)
 =\frac{b+1}{(2\sigma+b+1)(2\sigma+b-1)}
 \ge\frac1{z+h}-\frac{2s}{(z+h)^2}.
\]
This also holds when the lower index is zero.
The probability of no merger by step $m$ is at most
\[
 \exp\left(-\sum_{h<m}\frac1{z+h}
                     +2s\sum_{h<m}\frac1{(z+h)^2}\right)
 \le e^2\frac z{z+m}.
\]
Here the first sum is at least $\log((z+m)/z)$, and the second
multiplied by $2s$ is at most
$2s(z^{-1}+z^{-2})\le2$.
The total variation distance of the two distributions is no greater
than the nonmerger probability. Their $\ell^1$ distance is twice
total variation, giving \eqref{inf:coupling-bound}.
Every subsequent positive stencil of sum one is an $\ell^1$
contraction, so it preserves this estimate.
\end{proof}

\begin{lemma}[Clamped multiplication gain]\label{inf:clamped-smoothing}
For integers $c,h\ge0$ put $T=S_ct^h$ and $D=2c+2h$.
On clamped sources, that is on $\Gc_p$,
\begin{equation}\label{inf:clamped-gain}
 \nrm{M_T\Euler K}_{\Gc_p\to\Cc}
 \le C\frac{2^c\wt(D)}{p+D}.
\end{equation}
The constant is independent of $c,h$, the source indices, $p$ and
the admissible split.
\end{lemma}
\begin{proof}
Use a source $\Psi_{j,s}$ with $s\ge1$, and put $B=2j$,
$\beta=B+p-1$, $d_s=\beta+2s$.
For an angular output of degree $L=2k$ set
\[
 n_{\downarrow}=h+(2c+B-L)/2,\qquad
 n_{\uparrow}=h+(2c-B+L)/2.
\]
There are $n_{\downarrow}$ powers of $t$. First raise the source
parameter by $n_{\uparrow}$, then lower it with those
$n_{\downarrow}$ powers; the net change is
$n_{\uparrow}-n_{\downarrow}=L-B$.
The triangle support guarantees $n_{\downarrow},n_{\uparrow}\ge0$.
If $D\le8(B+p)$, the ordinary multiplier bound and the column
bound $C/d_s$ for $\Euler K$ prove the claim, since
$p+D\le9(B+p)\le9(d_s+1)$.

Suppose $D>8(B+p)$. Then $n_{\uparrow}\ge D/4$.
If $2c<D/2$, one has $h>D/4$ and $n_{\uparrow}\ge h$.
If $2c\ge D/2$, the triangle inequality gives
$n_{\uparrow}\ge h+2c-B>D/4$.
Since $s\ge1$, the Poisson column has the zero-sum form
\[
 k_s=\frac{J_{s-1}^\beta-J_s^\beta}{4d_s(d_s+1)}
       +\frac{J_{s+1}^\beta-J_s^\beta}{4(d_s+1)(d_s+2)}.
\]
After $n_{\uparrow}$ raises, Lemma~\ref{inf:raising-coupling} bounds its
coefficient sum by $C/(d_s(d_s+n_{\uparrow}))$.
The term $2tk_s'$ begins with the adjacent difference in
\eqref{inf:poisson-derivative}, in parameter $\beta+1$.
Raise it by $n_{\uparrow}$ and then lower with $n_{\downarrow}+1$
powers, including its extra $t$. Its coefficient sum is at most
$C/(d_s+n_{\uparrow})$.
Since $B\le d_s$, the combination $Bk_s+2tk_s'$ is bounded by
$C/(d_s+n_{\uparrow})\le C/(p+D)$.

All output degrees are at most the input degree plus $D+2$.
The extra factor $\wt(2)$ is bounded by two, and angular weights
are at most $2^{c+j}$. Summing the positive angular coefficients
in \eqref{inf:gasper-product}, and using
\eqref{inf:weight-properties}, proves \eqref{inf:clamped-gain}.
\end{proof}

\subsection{Gradient estimates with one shape moment}

\begin{lemma}[Gradient profiles on every source layer]
\label{inf:gradient-profiles}
Let $A=2i>0$, $B=2j$, $\ell_{\mathrm{out}}=2k$, and $h=(A+B-\ell_{\mathrm{out}})/2$ for an
angular product output in \eqref{inf:gasper-product}. Put
$Z_{\mathrm{grad}}=h+\ell_{\mathrm{out}}+p-1$.
For $\KD\Psi_{j,s}=S_jk_s$ the two output profiles
of $\nabla S_i\cdot\nabla\KD$ and
$\nabla S_i\cdot\nabla\Euler\KD$, with the coefficient
$c_{ij}^k$ suppressed, are respectively
\begin{align}
 &2hZ_{\mathrm{grad}}t^{h-1}k_s+2At^hk'_s,\label{inf:gradient-profile-one}\\
 &At^hJ_s^\beta+
 (4hZ_{\mathrm{grad}}-2AB-4A(p-1))t^hk'_s
                  +2BhZ_{\mathrm{grad}}t^{h-1}k_s.\label{inf:gradient-profile-two}
\end{align}
The term with $h=0$ and a negative power of $t$ is absent.
Consequently, including $s=0$,
\begin{align}
 \nrm{\nabla S_i\cdot\nabla\KD}_{\Cc\to\Cc}
 &\le CA2^i\wt(A)/p,\label{inf:gradient-C-one}\\
 \nrm{\nabla S_i\cdot\nabla\Euler\KD}_{\Cc\to\Cc}
 &\le CA2^i\wt(A).\label{inf:gradient-C-two}
\end{align}
\end{lemma}
\begin{proof}
Since both solid factors are harmonic,
$2\nabla S_i\cdot\nabla S_j=\Delta(S_iS_j)$.
Use $S_iS_j=\sum c_{ij}^kS_kt^h$ and
\eqref{inf:profile-laplacian}; its coefficient is $4hZ_{\mathrm{grad}}$.
The derivative of the radial profile contributes $2At^hk_s'$,
proving \eqref{inf:gradient-profile-one}.
Apply that formula to $Bk_s+2tk_s'$.
The Poisson equation gives
$(Bk_s+2tk_s')'=J_s^\beta/2-(B+2p-2)k_s'$;
substitution proves \eqref{inf:gradient-profile-two}.
The triangle support gives $0\le h\le\min(A,B)$ and
\begin{equation}\label{inf:hZ-bound}
 hZ_{\mathrm{grad}}=h(A+B-h+p-1)\le A(B+p-1).
\end{equation}
For $A\le B$ the left side is increasing up to $h=A$.
For $A\ge B$ use its value at $h=B$ and
$B(A+p-1)\le A(B+p-1)$.

In every nonsaturated output the positive parameter conversions of
Lemma~\ref{inf:radial-stencils} apply separately to the two terms
in \eqref{inf:gradient-profile-one}.
When $\ell_{\mathrm{out}}<B$, the required lowerings are $B-\ell_{\mathrm{out}}$ for the $k_s$ term
and $B-\ell_{\mathrm{out}}+1$ for the derivative term, whose parameter is $\beta+1$.
The available powers are $h-1$ and $h$, respectively.
They suffice whenever $\ell_{\mathrm{out}}\ge B-A+2$; the degrees are even.
When $\ell_{\mathrm{out}}=B$, the $k_s$ term keeps its parameter.
The derivative term starts in parameter $\beta+1$ and needs one
lowering; one of its $h=i\ge1$ powers of $t$ supplies it.
When $\ell_{\mathrm{out}}>B$, raise the $k_s$ term by
$\ell_{\mathrm{out}}-B$ and the derivative term by
$\ell_{\mathrm{out}}-B-1\ge1$, followed by their nonnegative
remaining powers of $t$. The case $h=0$ omits the first term.
Thus the only remaining case is $\ell_{\mathrm{out}}=B-A$, $h=A$, where the two
terms combine exactly to
\[
 2At^{A-1}(\beta k_s+tk_s').
\]
Use \eqref{inf:poisson-lowered} before the remaining $A-1$
lowerings. This supplies the one missing lowering step.
The separate harmonic-column identity in Lemma~\ref{inf:poisson-columns}
proves the same statement at $s=0$.

The absolute Poisson column sums are $O((B+p+2s)^{-2})$, and
the two first-derivative columns have sums $O((B+p+2s)^{-1})$.
Together with \eqref{inf:hZ-bound}, these give a per-output bound
$CA/(B+p)$ for \eqref{inf:gradient-profile-one}.
In \eqref{inf:gradient-profile-two} the three terms have bounds
$CA$, $4CA$, and $2CA$ after positive conversion.
In the saturated case use their combined expression
\[
 At^AJ_s^\beta+2ABt^{A-1}(\beta k_s+tk_s'),
\]
which gives the same bound, also at $s=0$.
In that layer the coefficient norm of $k_0$ is
$1/(2b_0(b_0+1))$, which supplies the factor $1/(B+p)$.
All final degrees are at most input degree plus $A$.
The positive angular coefficients sum to one and
$2^k\le2^{i+j}$. Applying \eqref{inf:weight-properties} proves
\eqref{inf:gradient-C-one}--\eqref{inf:gradient-C-two}.
\end{proof}

\begin{proposition}[Grouped shape derivatives]\label{inf:grouped-shapes}
For $H_i=\HP h_i$, the following bounds hold, with absolute
constants:
\begin{align}
 \nrm{M_{\nabla H_1\cdot\nabla H_2}}_{\Cc\to\Cc}
 &\le Cp\nrm{h_1}_{\Ec_1}\nrm{h_2}_{\Ec_1},
 \label{inf:gradient-multiplier}\\
 \nrm{M_{\Euler(\Euler-1)H_1}\Euler K}_{\Gc_p\to\Cc}
 &\le C\nrm{h_1}_{\Ec_1},\label{inf:second-shape-clamped}\\
 \nrm{M_{\nabla H_1\cdot\nabla\Euler H_2}\Euler K}_{\Gc_p\to\Cc}
 &\le C\nrm{h_1}_{\Ec_1}\nrm{h_2}_{\Ec_1}.
 \label{inf:mixed-shape-clamped}
\end{align}
The last two bounds are on clamped sources in $\Gc_p$.
\end{proposition}
\begin{proof}
For a pair of harmonics of degrees $A=2i,B=2j$, the product of
gradients has output coefficient $2hZ_{\mathrm{grad}}$ and solid profile
$S_kt^{h-1}$. It is zero if $h=0$; otherwise its total degree is
$A+B-2$. From \eqref{inf:hZ-bound},
$2hZ_{\mathrm{grad}}\le Cp(1+i)(1+j)$.
The module bound and positive angular summation prove
\eqref{inf:gradient-multiplier}.
For the second inequality apply Lemma~\ref{inf:clamped-smoothing}
to the harmonic degree $A$, with coefficient $A(A-1)$.
The quotient $A(A-1)/(p+A)$ is at most $A$, leaving one shape
moment.
For the last inequality its pairwise coefficient is $2BhZ_{\mathrm{grad}}$.
Apply the same clamped gain at degree $D=A+B-2$.
For $h\ge1$, $Z_{\mathrm{grad}}=A+B-h+p-1\le p+A+B-2$, so
\[
 \frac{2BhZ_{\mathrm{grad}}}{p+A+B-2}\le2B\min(A,B)\le2AB.
\]
The grading and angular weights are bounded by the product of the
two shape weights. Summation gives \eqref{inf:mixed-shape-clamped}.
\end{proof}

\subsection{Double-trace lifts and source moments}

\begin{lemma}[Exact double-trace lift]\label{inf:double-lift}
For each angular index define
\begin{equation}\label{inf:lift-columns}
 L_j^D=S_j[1-j(t-1)],\qquad L_j^N=\tfrac12S_j(t-1).
\end{equation}
Their Dirichlet and radial normal traces are $(G_j,0)$ and
$(0,G_j)$, respectively. Their coefficient bounds and Laplacians
are
\begin{align}
 \nrm{L_j^D}_\Cc&\le3\,2^j\wt(2j+2),&
 \nrm{L_j^N}_\Cc&\le\frac{2^j\wt(2j+2)}{2j+p+1},
 \label{inf:lift-column-bounds}\\
 \Delta L_j^D&=-2j(4j+2p)S_j,&
 \Delta L_j^N&=(4j+2p)S_j.\label{inf:lift-laplacians}
\end{align}
For the defects of Corollary~\ref{inf:centre-coefficient-defect},
let $f_N=h_N/R$, $q_N=1+\HP f_N$ and
$d_N=q_N+\Euler q_N$. Define
\begin{equation}\label{inf:reference-defect}
 N_N^{\mathrm{ref}}=R(d_N|_{\partial B})N_N^{\mathrm{phys}},
 \qquad
 \ell_N=\sum_j(D_{N,j}L_j^D+N_{N,j}^{\mathrm{ref}}L_j^N).
\end{equation}
For every fixed moment $k$ and $p\ge P_{N,k}$,
\begin{equation}\label{inf:lift-estimates}
 \nrm{\ell_N}_{\Cc^{[k]}}\le C_{N,k}\theta^2p^{-N},\qquad
 \nrm{\Delta\ell_N}_{\Cc^{[k]}}\le C_{N,k}\theta^2p^{2-N}.
\end{equation}
\end{lemma}
\begin{proof}
Differentiating \eqref{inf:lift-columns} at $t=1$ gives their
traces. The identity
$t-1=(J_1^{\beta_j}-J_0^{\beta_j})/(\beta_j+2)$
places both lifts in radial layers zero and one and gives
\eqref{inf:lift-column-bounds}. Formula
\eqref{inf:profile-laplacian} proves \eqref{inf:lift-laplacians}.
The reference normal trace in \eqref{inf:reference-defect} follows
by differentiating the physical composition $u_N(Rq_N(y)y)$ in
the radial direction on the sphere: $D(q_Ny)\nu=d_N\nu$.

The multiplier $d_N|_{\partial B}$ is a fixed-degree polynomial
with uniformly bounded coefficients at any fixed angular radius.
Thus \eqref{inf:centre-coeff-defect-bound} bounds its product with
the physical defect, losing only the factor $R=O(p)$.
The factor $(2j+p+1)^{-1}$ in the normal lift cancels that factor.
Choose a slightly larger fixed angular radius in
\eqref{inf:centre-coeff-defect-bound} to absorb
$\wt(2j+2)\le e^{(2j+2)/p}$ and the prescribed angular moments.
This proves the first estimate in \eqref{inf:lift-estimates}.
The Laplacians in \eqref{inf:lift-laplacians}, including the extra
$p$ in the reference normal trace, cost at most $C p^2$ after
absorbing fixed powers of $j$ by the same radius choice. This proves
the second estimate. The sums converge in all stipulated moment
norms and their traces agree with the classical lifts by the same
coefficient identities and smooth convergence.
\end{proof}

\begin{lemma}[Derivative columns and moments]\label{inf:derivative-moments}
For $\beta\ge0$, $s\ge1$, and $0\le k<s$,
\begin{equation}\label{inf:derivative-column}
 \partial_tJ_s^\beta=\sum_{k<s}d_{s,k}J_k^\beta,
 \quad d_{s,k}=(\beta+2k+1)
 \left(1-(-1)^{s+k}\frac{\binom{k+\beta}{k}}{\binom{s+\beta}{s}}\right)
 \ge0.
\end{equation}
For every $0\le m\le s$ the endpoint derivative is
\begin{equation}\label{inf:Jacobi-endpoint-derivatives}
 (J_s^\beta)^{(m)}(1)
 =\frac{(s)_{\mathrm{fall},m}(s+\beta+1)_m}{m!}.
\end{equation}
In particular, the first derivative column has sum
$s(s+\beta+1)$ and the second has sum
$s(s-1)(s+\beta+1)(s+\beta+2)/2$.
Consequently
\begin{equation}\label{inf:Euler-moment}
 \nrm{\Euler u}_{\Cc^{[k]}}\le Cp\nrm u_{\Cc^{[k+2]}}.
\end{equation}
The Laplacian has positive radial columns of sum at most
$C(p+2j+2s)^4$ and does not raise total degree. For each fixed $k\ge0$,
\begin{equation}\label{inf:Delta-moment}
 \nrm{\Delta u}_{\Cc^{[k]}}
 \le C_k\{p^4\nrm u_{\Cc^{[k]}}+\nrm u_{\Cc^{[k+4]}}\}.
\end{equation}
\end{lemma}
\begin{proof}
For $\beta>0$, integration by parts gives, since $k<s$,
\[
 \int_0^1t^\beta(J_s^\beta)'J_k^\beta\,dt
 =1-\beta J_k^\beta(0)\int_0^1t^{\beta-1}J_s^\beta\,dt.
\]
To obtain this expression, subtract $J_k^\beta(0)$ in the
integrated derivative term; its remaining factor is $t$ times a
polynomial of degree below $s$, hence orthogonal to $J_s^\beta$.
Rodrigues' formula and $s$ integrations by parts give the remaining
integral as $(-1)^sB(\beta,s+1)$.
Divide by the squared norm $(\beta+2k+1)^{-1}$ and use the endpoint
value at zero to obtain \eqref{inf:derivative-column}.
The formula at $\beta=0$ is its finite continuous limit.
Monotonicity in $k$ of $\binom{k+\beta}{k}$ proves nonnegativity.
Every further differentiation preserves coefficient positivity.
For the endpoint values, write $J_s^\beta(1+z)=\sum_m\mathfrak d_mz^m$
and use its differential equation
\[
 t(1-t)J''+\{\beta+1-(\beta+2)t\}J'
                      +s(s+\beta+1)J=0.
\]
Its coefficient of $z^m$ gives
$\mathfrak d_{m+1}=(s-m)(s+\beta+1+m)\mathfrak d_m/(m+1)^2$,
with $\mathfrak d_0=1$.
Iteration and multiplication by $m!$ prove
\eqref{inf:Jacobi-endpoint-derivatives}. Evaluating the positive
columns at one now gives their sums.

On one solid profile $\Euler=B+2t\partial_t$.
The positive $t$ stencil gives column sum
$B+2s(s+\beta+1)\le Cp(1+2j+2s)^2$ and no increase in total degree.
This proves \eqref{inf:Euler-moment}.
The profile Laplacian is $4(t\partial_t^2+(\beta+1)\partial_t)$,
so its column sum is
$4((J_s^\beta)''(1)+(\beta+1)(J_s^\beta)'(1))$,
bounded by $C(p+2j+2s)^4$.
Each output degree is at most the input degree $D=2j+2s$, so
its moment weight $(1+D_{\mathrm{out}})^k$ is at most $(1+D)^k$.
Multiplying the column bound by this weight and using
$(p+D)^4\le8(p^4+D^4)$ gives \eqref{inf:Delta-moment} after
absolute summation.
\end{proof}
